\documentclass[10pt,a4paper]{amsart}
\usepackage{amsmath,amssymb,amsthm,graphicx}
% Retained minimal preamble for the Stacks Project chapter sites-cohomology.tex (Cohomology on Sites).
% Verbatim excerpt of the chapter preamble: Xy-pic setup (native preamble lines 14-18) and the native
% theorem declarations with the mathematical macros (native preamble lines 154-217). Copied mechanically;
% nothing added, removed or redefined. The wrapper class and the amsmath/amssymb/amsthm/graphicx base
% packages are supplied by the assembler, so the class/externaldocument part of the native preamble is omitted.
% For commutative diagrams we use Xy-pic
\usepackage[all]{xy}

% We use 2cell for 2-commutative diagrams.
\xyoption{2cell}
% Theorem environments.
%
\theoremstyle{plain}
\newtheorem{theorem}[subsection]{Theorem}
\newtheorem{proposition}[subsection]{Proposition}
\newtheorem{lemma}[subsection]{Lemma}

\theoremstyle{definition}
\newtheorem{definition}[subsection]{Definition}
\newtheorem{example}[subsection]{Example}
\newtheorem{exercise}[subsection]{Exercise}
\newtheorem{situation}[subsection]{Situation}

\theoremstyle{remark}
\newtheorem{remark}[subsection]{Remark}
\newtheorem{remarks}[subsection]{Remarks}

\numberwithin{equation}{subsection}

% Macros
%
\def\lim{\mathop{\mathrm{lim}}\nolimits}
\def\colim{\mathop{\mathrm{colim}}\nolimits}
\def\Spec{\mathop{\mathrm{Spec}}}
\def\Hom{\mathop{\mathrm{Hom}}\nolimits}
\def\Ext{\mathop{\mathrm{Ext}}\nolimits}
\def\SheafHom{\mathop{\mathcal{H}\!\mathit{om}}\nolimits}
\def\SheafExt{\mathop{\mathcal{E}\!\mathit{xt}}\nolimits}
\def\Sch{\mathit{Sch}}
\def\Mor{\mathop{\mathrm{Mor}}\nolimits}
\def\Ob{\mathop{\mathrm{Ob}}\nolimits}
\def\Sh{\mathop{\mathit{Sh}}\nolimits}
\def\NL{\mathop{N\!L}\nolimits}
\def\CH{\mathop{\mathrm{CH}}\nolimits}
\def\proetale{{pro\text{-}\acute{e}tale}}
\def\etale{{\acute{e}tale}}
\def\QCoh{\mathit{QCoh}}
\def\Ker{\mathop{\mathrm{Ker}}}
\def\Im{\mathop{\mathrm{Im}}}
\def\Coker{\mathop{\mathrm{Coker}}}
\def\Coim{\mathop{\mathrm{Coim}}}

% Boxtimes
%
\DeclareMathSymbol{\boxtimes}{\mathbin}{AMSa}{"02}

%
% Macros for moduli stacks/spaces
%
\def\QCohstack{\mathcal{QC}\!\mathit{oh}}
\def\Cohstack{\mathcal{C}\!\mathit{oh}}
\def\Spacesstack{\mathcal{S}\!\mathit{paces}}
\def\Quotfunctor{\mathrm{Quot}}
\def\Hilbfunctor{\mathrm{Hilb}}
\def\Curvesstack{\mathcal{C}\!\mathit{urves}}
\def\Polarizedstack{\mathcal{P}\!\mathit{olarized}}
\def\Complexesstack{\mathcal{C}\!\mathit{omplexes}}
% \Pic is the operator that assigns to X its picard group, usage \Pic(X)
% \Picardstack_{X/B} denotes the Picard stack of X over B
% \Picardfunctor_{X/B} denotes the Picard functor of X over B
\def\Pic{\mathop{\mathrm{Pic}}\nolimits}
\def\Picardstack{\mathcal{P}\!\mathit{ic}}
\def\Picardfunctor{\mathrm{Pic}}
\def\Deformationcategory{\mathcal{D}\!\mathit{ef}}
\usepackage{cleveref}
\crefname{originaldefinition}{original definition}{original definitions}
\Crefname{originaldefinition}{Original definition}{Original definitions}
\makeatletter
\newlabel{categories-definition-dual}{{0FFP}{source}{}{sourcepointer-1}{}}
\newlabel{categories-definition-dual@cref}{{[originaldefinition][0][]0FFP}{[1][0][]source}{}{}{}}
\makeatother
\crefname{originalremark}{original remark}{original remarks}
\Crefname{originalremark}{Original remark}{Original remarks}
\makeatletter
\newlabel{categories-remark-left-dual-adjoint}{{0FFR}{source}{}{sourcepointer-2}{}}
\newlabel{categories-remark-left-dual-adjoint@cref}{{[originalremark][0][]0FFR}{[1][0][]source}{}{}{}}
\makeatother
\crefname{originalreference}{original reference}{original references}
\Crefname{originalreference}{Original reference}{Original references}
\makeatletter
\newlabel{example-dual}{{0FNV}{source}{}{sourcepointer-3}{}}
\newlabel{example-dual@cref}{{[originalreference][0][]0FNV}{[1][0][]source}{}{}{}}
\makeatother
\crefname{originallemma}{original lemma}{original lemmas}
\Crefname{originallemma}{Original lemma}{Original lemmas}
\makeatletter
\newlabel{sites-modules-lemma-left-dual-module}{{0FNZ}{source}{}{sourcepointer-4}{}}
\newlabel{sites-modules-lemma-left-dual-module@cref}{{[originallemma][0][]0FNZ}{[1][0][]source}{}{}{}}
\makeatother
\title{A complex of modules on a ringed site with a left dual is locally strictly perfect and its dual is the termwise one}
\author{The Stacks Project Authors}
\date{stacks-project:sites-cohomology.tex:local-snapshot-20261011; GNU FDL 1.2 or later}
\begin{document}
\maketitle
\noindent\textit{Editorial note.} All delivered mathematical text is the author's own, copied from the pinned Stacks Project chapter file sites-cohomology.tex; no statement or proof step was rewritten, reordered or removed.\par\medskip
\noindent\textit{Editorial note.} The counted proof is the author's complete argument: the proof is written out in full: the reduction through uniqueness of left duals, the decomposition of eta and epsilon with the resulting termwise duality, and the two localisation steps producing finitely many nonzero terms and terms that are locally direct summands of finite free modules.\par\medskip
\noindent\textit{Editorial note.} The label lemma-left-dual-complex names the counted statement, and the proof environment immediately following it in the source is delivered verbatim as its proof.\par\medskip
\noindent\textit{Editorial note.} The counted statement is carried by the author's own lemma environment in the chapter source.\par\medskip
\noindent\textit{Editorial note.} The statement and proof are copied from the pinned chapter file; the four cited inputs are kept as external pointers with their canonical Stacks tags, including the example reference whose kind is recorded as section as the assembler requires for example-labelled pointers.\par\medskip
\noindent\textit{Editorial note.} The preamble is the chapter's own retained minimal preamble (sources/stacks-sites-cohomology/retained-preamble.tex), a verbatim excerpt of the native preamble: the Xy-pic setup and the native theorem declarations with the mathematical macros.\par\medskip
\noindent\textit{Allowed standard prerequisites (editorial).}\begin{itemize}
\item Left dual in a monoidal category and uniqueness of duals. Categories, the definition of a left dual through the unit and counit maps, and the remark that left duals are unique up to canonical isomorphism. Supplies the data G\textasciicircum{}bullet, eta, epsilon and the reduction of the final statement to the local structure assertion.
\item The termwise left dual of a module. Modules on Sites, the lemma describing left duals of O-modules and of direct summands of finite free modules. Supplies the conclusion that each F\textasciicircum{}n is locally a direct summand of a finite free module once finitely many terms are nonzero.
\item The explicit left dual constructed on a strictly perfect complex. Same chapter, the example constructing the left dual of a strictly perfect complex. Supplies the identification of the abstract left dual with the termwise one asserted at the end.
\end{itemize}\medskip
\ifdefined\hypertarget\hypertarget{sourcepointer-1}{}\fi
\noindent\textit{Source reference.} Original definition 0FFP: sources/stacks-chow/tags.txt, source lines 19294--19294. Prerequisite pointer invoked at proof line(s) 13005; it names a standard fact of a sibling chapter that the counted argument uses, and every step of the counted proof built on it is retained verbatim in the delivered body.\par\medskip
\ifdefined\hypertarget\hypertarget{sourcepointer-2}{}\fi
\noindent\textit{Source reference.} Original remark 0FFR: sources/stacks-chow/tags.txt, source lines 19296--19296. Prerequisite pointer invoked at proof line(s) 13014; it names a standard fact of a sibling chapter that the counted argument uses, and every step of the counted proof built on it is retained verbatim in the delivered body.\par\medskip
\ifdefined\hypertarget\hypertarget{sourcepointer-3}{}\fi
\noindent\textit{Source reference.} Original reference 0FNV: sources/stacks-sites-cohomology/pinned-original-source.tex, source lines 12960--12960. Prerequisite pointer invoked at proof line(s) 13009; it names a standard fact of the same chapter that the counted argument uses, and every step of the counted proof built on it is retained verbatim in the delivered body.\par\medskip
\ifdefined\hypertarget\hypertarget{sourcepointer-4}{}\fi
\noindent\textit{Source reference.} Original lemma 0FNZ: sources/stacks-chow/tags.txt, source lines 19584--19584. Prerequisite pointer invoked at proof line(s) 13034; it names a standard fact of a sibling chapter that the counted argument uses, and every step of the counted proof built on it is retained verbatim in the delivered body.\par\medskip
\begin{lemma}
\label{lemma-left-dual-complex}
Let $(\mathcal{C}, \mathcal{O})$ be a ringed site. Let $\mathcal{F}^\bullet$
be a complex of $\mathcal{O}$-modules. If $\mathcal{F}^\bullet$
has a left dual in the monoidal category of complexes of
$\mathcal{O}$-modules
(Categories, Definition \ref{categories-definition-dual})
then for every object $U$ of $\mathcal{C}$ there exists a
covering $\{U_i \to U\}$ such that $\mathcal{F}^\bullet|_{U_i}$
is strictly perfect and the left dual is as constructed in
Example \ref{example-dual}.
\end{lemma}

\begin{proof}
By uniqueness of left duals
(Categories, Remark \ref{categories-remark-left-dual-adjoint})
we get the final statement provided we show that $\mathcal{F}^\bullet$
is as stated. Let $\mathcal{G}^\bullet, \eta, \epsilon$ be a left dual.
Write $\eta = \sum \eta_n$ and $\epsilon = \sum \epsilon_n$
where $\eta_n : \mathcal{O} \to
\mathcal{F}^n \otimes_\mathcal{O} \mathcal{G}^{-n}$
and
$\epsilon_n : \mathcal{G}^{-n} \otimes_\mathcal{O} \mathcal{F}^n
\to \mathcal{O}$. Since
$(1 \otimes \epsilon) \circ (\eta \otimes 1) = \text{id}_{\mathcal{F}^\bullet}$
and
$(\epsilon \otimes 1) \circ (1 \otimes \eta) = \text{id}_{\mathcal{G}^\bullet}$
by Categories, Definition \ref{categories-definition-dual} we see immediately
that we have
$(1 \otimes \epsilon_n) \circ (\eta_n \otimes 1) = \text{id}_{\mathcal{F}^n}$
and
$(\epsilon_n \otimes 1) \circ (1 \otimes \eta_n) =
\text{id}_{\mathcal{G}^{-n}}$.
In other words, we see that $\mathcal{G}^{-n}$ is a left dual of
$\mathcal{F}^n$ and we see that
Modules on Sites, Lemma \ref{sites-modules-lemma-left-dual-module}
applies to each $\mathcal{F}^n$. Let $U$ be an object of $\mathcal{C}$.
There exists a covering $\{U_i \to U\}$ such that for every
$i$ only a finite number of $\eta_n|_{U_i}$ are nonzero.
Thus after replacing $U$ by $U_i$ we may assume only a finite
number of $\eta_n|_U$ are nonzero and by the lemma cited
this implies only a finite number of $\mathcal{F}^n|_U$ are 
nonzero. Using the lemma again we can then find a covering
$\{U_i \to U\}$ such that each
$\mathcal{F}^n|_{U_i}$ is a direct summand of a finite
free $\mathcal{O}$-module and the proof is complete.
\end{proof}

\end{document}
